\documentclass[11pt,a4paper]{article}
\usepackage[margin=1in]{geometry}
\usepackage{amsmath,amssymb,amsthm}
\usepackage{algorithm}
\usepackage{algpseudocode}
\usepackage{booktabs}
\usepackage{hyperref}

\theoremstyle{definition}
\newtheorem{definition}{Definition}[section]
\newtheorem{theorem}{Theorem}[section]
\newtheorem{lemma}[theorem]{Lemma}
\newtheorem{remark}{Remark}[section]

\title{\textbf{Rigorous Mathematical Specification, Sign Momentum Dynamics, and Convergence of the Lion Optimizer}}
\author{Team Scaibu \\ \small Email: \texttt{jaydeep.9664920749@gmail.com} \\ \small Architecture and Core Engine Team}
\date{October 2026}

\begin{document}
\maketitle

\begin{abstract}
This document presents the rigorous mathematical foundation and algorithmic specification of the EvoLved Sign Momentum (Lion) optimizer. Discovered via symbolic program search, Lion replaces coordinate-wise second-moment division with the signed projection of an interpolated momentum vector. This halves auxiliary optimizer memory compared to Adam while maintaining uniform update scales ($\|\Delta \theta\|_\infty = \eta$). We formalize the dual-momentum tracking dynamics, provide complete algorithmic pseudo-code matching the reference implementation, prove convergence under non-convex smooth objectives, derive computational complexities, demonstrate a step-by-step numerical calculation, and address floating-point considerations.
\end{abstract}

\section{Notation and Mathematical Preliminaries}

Let $\theta \in \mathbb{R}^P$ denote the network parameter vector trained across discrete steps $t \in \{1, \dots, T\}$.

\begin{itemize}
    \item $g_t \triangleq \nabla_\theta \mathcal{L}(\theta_t) \in \mathbb{R}^P$: Stochastic gradient vector at iterate $\theta_t$.
    \item $m_t \in \mathbb{R}^P$: Exponential moving average momentum buffer.
    \item $c_t \in \mathbb{R}^P$: Interpolated momentum update vector at step $t$.
    \item $\operatorname{sign}(z)$: Coordinate-wise signum function, defined as $+1$ if $z > 0$, $-1$ if $z < 0$, and $0$ if $z = 0$.
    \item $\beta_1 \in [0, 1)$: Momentum interpolation coefficient for update direction (default $0.9$).
    \item $\beta_2 \in [0, 1)$: Momentum decay coefficient for tracking buffer (default $0.99$).
    \item $\lambda \ge 0$: Decoupled weight decay coefficient (default $0.1$).
    \item $\eta > 0$: Base learning rate (typically $3\times$ to $10\times$ smaller than Adam, e.g., $10^{-4}$).
\end{itemize}

\begin{table}[h]
\centering
\caption{Summary of Mathematical Symbols for Lion}
\begin{tabular}{lll}
\toprule
\textbf{Symbol} & \textbf{Type / Domain} & \textbf{Description} \\
\midrule
$\theta_t$ & $\mathbb{R}^P$ & Parameter iterate at step $t$ \\
$g_t$ & $\mathbb{R}^P$ & Stochastic gradient vector $\nabla_\theta \mathcal{L}(\theta_t)$ \\
$m_t$ & $\mathbb{R}^P$ & Tracking momentum buffer \\
$c_t$ & $\mathbb{R}^P$ & Interpolated update momentum $\beta_1 m_{t-1} + (1 - \beta_1) g_t$ \\
$\beta_1, \beta_2$ & $[0, 1)$ & Interpolation and tracking decay factors ($0.9, 0.99$) \\
$\lambda$ & $\mathbb{R}_{\ge 0}$ & Decoupled weight decay ($0.1$) \\
$\eta$ & $\mathbb{R}_{> 0}$ & Base learning rate ($10^{-4}$) \\
\bottomrule
\end{tabular}
\end{table}

\section{Problem Definition and Core Concepts}

\subsection{Intuition and Motivation}
Adam tracks both first moments ($m \in \mathbb{R}^P$) and second moments ($v \in \mathbb{R}^P$), computing coordinate updates $\Delta \theta_i \propto \frac{\hat{m}_i}{\sqrt{\hat{v}_i}}$. The Lion optimizer simplifies this mechanism: instead of estimating individual coordinate variances, it takes the sign of an interpolated momentum vector $c_t = \beta_1 m_{t-1} + (1 - \beta_1) g_t$. Because every coordinate update has identical magnitude $\pm \eta$, Lion acts as a natural regularizer, bounding update step norms in the $L_\infty$ ball ($\|\Delta \theta\|_\infty \le \eta$) and cutting optimizer state memory by $50\%$ by eliminating the second-moment tensor $v$.

\subsection{Formal Mathematical Definition}
\begin{definition}[Lion Optimizer Update Rule]
Given initial parameter vector $\theta_0 \in \mathbb{R}^P$, initial momentum $m_0 = \mathbf{0} \in \mathbb{R}^P$, and hyperparameters $\eta > 0, \lambda \ge 0, \beta_1, \beta_2 \in [0, 1)$, the recurrence relations for step $t \ge 1$ and coordinate $i \in \{1, \dots, P\}$ are:
\begin{align}
c_{t, i} &= \beta_1 m_{t-1, i} + (1 - \beta_1) g_{t, i} \label{eq:lion_interp} \\
\theta_{t, i} &= \theta_{t-1, i} (1 - \eta \lambda) - \eta \operatorname{sign}(c_{t, i}) \label{eq:lion_update} \\
m_{t, i} &= \beta_2 m_{t-1, i} + (1 - \beta_2) g_{t, i} \label{eq:lion_buffer}
\end{align}
\end{definition}

\section{Algorithmic Specification}

The exact operational algorithm implemented in \texttt{NnAlgoLionOptimizer} is shown below.

\begin{algorithm}[H]
\caption{Lion (EvoLved Sign Momentum) Step Execution}
\begin{algorithmic}[1]
\Require Parameter vector $\theta \in \mathbb{R}^P$, gradient vector $g \in \mathbb{R}^P$, momentum buffer $m \in \mathbb{R}^P$
\Require Step size $\eta > 0$, decays $\beta_1, \beta_2 \in [0, 1)$, weight decay $\lambda \ge 0$
\Ensure Updated parameters $\theta_{\text{next}} \in \mathbb{R}^P$, updated momentum buffer $m_{\text{next}} \in \mathbb{R}^P$
\State $P \gets \text{length}(\theta)$
\If{$P = 0$ \textbf{or} $\text{length}(g) \ne P$ \textbf{or} $\text{length}(m) \ne P$}
    \State \textbf{throw} PreconditionFailureException(``Buffer dimensions mismatch or empty'')
\EndIf
\If{$\eta \le 0$ \textbf{or} $\lambda < 0$ \textbf{or} $\beta_1 < 0$ \textbf{or} $\beta_1 \ge 1$ \textbf{or} $\beta_2 < 0$ \textbf{or} $\beta_2 \ge 1$}
    \State \textbf{throw} PreconditionFailureException(``Hyperparameter domain violation'')
\EndIf
\State $\theta_{\text{next}} \gets \text{new Array of size } P$
\State $m_{\text{next}} \gets \text{new Array of size } P$
\For{$i \gets 1 \textbf{ to } P$}
    \State $c_i \gets \beta_1 \cdot m[i] + (1.0 - \beta_1) \cdot g[i]$
    \State $s_i \gets \begin{cases} 1.0 & \text{if } c_i > 0.0 \\ -1.0 & \text{if } c_i < 0.0 \\ 0.0 & \text{otherwise} \end{cases}$
    \State $\theta_{\text{next}}[i] \gets \theta[i] \cdot (1.0 - \eta \cdot \lambda) - \eta \cdot s_i$
    \State $m_{\text{next}}[i] \gets \beta_2 \cdot m[i] + (1.0 - \beta_2) \cdot g[i]$
\EndFor
\State \Return $\{\text{updated\_parameters}: \theta_{\text{next}}, \text{updated\_exp\_avg}: m_{\text{next}}\}$
\end{algorithmic}
\end{algorithm}

\section{Theoretical Guarantees and Convergence Analysis}

\begin{theorem}[Convergence of Sign-SGD / Lion under Smooth Objectives]
Let $\mathcal{L}: \mathbb{R}^P \to \mathbb{R}$ be $L$-smooth. Assume gradients have bounded coordinate noise $\mathbb{E}[|g_{t, i} - \nabla_i \mathcal{L}(\theta_t)|^2] \le \sigma^2$. Under an appropriately chosen step-size schedule $\eta_t = \frac{\eta_0}{\sqrt{t}}$ and momentum decay, the expected $L_1$-norm of the objective gradient converges as:
\begin{equation}
\min_{t=1, \dots, T} \mathbb{E}\left[\frac{1}{P} \|\nabla \mathcal{L}(\theta_t)\|_1\right] \le \mathcal{O}\left(\frac{1}{T^{1/4}}\right)
\end{equation}
\end{theorem}

\begin{proof}
Let $\Delta \theta_t = -\eta \operatorname{sign}(c_t)$. Under $L$-smoothness:
\begin{equation}
\mathcal{L}(\theta_{t+1}) \le \mathcal{L}(\theta_t) - \eta \langle \nabla \mathcal{L}(\theta_t), \operatorname{sign}(c_t) \rangle + \frac{L \eta^2}{2} \|\operatorname{sign}(c_t)\|_2^2
\end{equation}
Notice $\|\operatorname{sign}(c_t)\|_2^2 \le P$ strictly for any dimension. Taking expectation with respect to stochastic batch gradients:
\begin{equation}
\mathbb{E}[\mathcal{L}(\theta_{t+1})] \le \mathbb{E}[\mathcal{L}(\theta_t)] - \eta \mathbb{E}[\|\nabla \mathcal{L}(\theta_t)\|_1] + \mathcal{O}(\eta \sigma) + \frac{L \eta^2 P}{2}
\end{equation}
Summing over $t=1 \dots T$ and choosing $\eta = \mathcal{O}(1 / \sqrt{P T})$ bounds the average gradient $L_1$ norm by $\mathcal{O}(T^{-1/4})$.
\end{proof}

\subsection{Limitations}
\begin{itemize}
    \item \textbf{Small Batch Instability}: In extremely small batch settings (e.g., batch size $< 64$), noisy gradient signs cause jittery updates; Lion excels most in large-batch training regimes ($B \ge 512$).
\end{itemize}

\section{Complexity Analysis}

\subsection{Time Complexity}
\begin{itemize}
    \item \textbf{FLOPs per Parameter}: Each coordinate requires 2 multiply-adds (for $c$ and $m$), 1 comparison (for sign), and 1 multiply-subtract (for $\theta$).
    \item \textbf{Total Time}: $\Theta(P)$ FLOPs per step (significantly faster than Adam since no square root or division is performed).
\end{itemize}

\subsection{Space Complexity}
\begin{itemize}
    \item \textbf{Auxiliary Memory}: Exactly one persistent vector $m \in \mathbb{R}^P$ ($4P$ bytes in FP32 vs. $8P$ for Adam).
    \item \textbf{Total Auxiliary Space}: $\Theta(P)$ floats (50\% memory reduction relative to Adam).
\end{itemize}

\section{Worked Numerical Example}

Consider $P=2$, $\theta = [1.0, -1.0]^T$, $g = [0.5, -0.2]^T$, $m = [0.1, -0.05]^T$, $\eta = 0.001$, $\beta_1 = 0.9$, $\beta_2 = 0.99$, $\lambda = 0.1$.

\subsection{Coordinate 1 Computation}
\begin{enumerate}
    \item $c_1 = 0.9(0.1) + (1 - 0.9)(0.5) = 0.09 + 0.05 = 0.14$.
    \item $\operatorname{sign}(c_1) = +1.0$ (since $0.14 > 0$).
    \item Weight update:
    \begin{align}
    \theta_{\text{next}, 1} &= 1.0 \times (1.0 - 0.001 \times 0.1) - 0.001 \times (+1.0) \\
    &= 1.0 \times (0.9999) - 0.001 = 0.9999 - 0.001 = 0.998900
    \end{align}
    \item Momentum update:
    \begin{equation}
    m_{\text{next}, 1} = 0.99(0.1) + 0.01(0.5) = 0.099 + 0.005 = 0.1040
    \end{equation}
\end{enumerate}

\subsection{Coordinate 2 Computation}
\begin{enumerate}
    \item $c_2 = 0.9(-0.05) + 0.1(-0.2) = -0.045 - 0.020 = -0.065$.
    \item $\operatorname{sign}(c_2) = -1.0$.
    \item Weight update:
    \begin{align}
    \theta_{\text{next}, 2} &= -1.0 \times 0.9999 - 0.001 \times (-1.0) = -0.9999 + 0.001 = -0.998900
    \end{align}
    \item Momentum update:
    \begin{equation}
    m_{\text{next}, 2} = 0.99(-0.05) + 0.01(-0.2) = -0.0495 - 0.002 = -0.0515
    \end{equation}
\end{enumerate}

\section{Numerical Considerations and Edge Cases}

\begin{itemize}
    \item \textbf{Sign of Zero}: If $c_i = 0.0$, the sign evaluates to $0.0$, applying purely decoupled weight decay without gradient displacement.
    \item \textbf{Quantization and Compression}: Because Lion updates are purely in $\{-1, 0, +1\}$, gradient updates across distributed worker nodes can be trivially 1-bit quantized with zero accuracy loss.
\end{itemize}

\begin{thebibliography}{9}
\bibitem{chen2023} X.~Chen et al., ``Symbolic Discovery of Optimization Algorithms,'' in \emph{Advances in Neural Information Processing Systems (NeurIPS)}, vol.~36, 2023.
\bibitem{bernstein2018} J.~Bernstein et al., ``signSGD: Compressed optimisation for non-convex problems,'' in \emph{International Conference on Machine Learning (ICML)}, pp.~560--569, 2018.
\bibitem{kingma2015} D.~P.~Kingma and J.~Ba, ``Adam: A method for stochastic optimization,'' in \emph{International Conference on Learning Representations (ICLR)}, 2015.
\end{thebibliography}

\end{document}
